% Source PDF: source/普通高考/2023/2023全国甲文(广西,贵州,西藏,四川).pdf
% Page: 1; pdfimages -list found no embedded bitmap; source page rendered at 300 dpi.
% Grid evidence: tmp/redraw_flowchart_2023_national_paper_a_liberal/grid/page-1-grid100.jpg,
%   q6_block_grid50.jpg, and q6_flow_block2_grid25.jpg.
% Original comparison crop: tmp/redraw_flowchart_2023_national_paper_a_liberal/grid/q6_original_final.png,
%   cropped from the no-grid page render.
% Elements: rounded start/end nodes, input/output parallelograms, decision diamond,
%   three assignment rectangles, directed connectors and branch labels.
% Complete node -> directed-edge list from the source:
%   start -> input -> test; test --是--> stepa -> stepb -> stepk -> test;
%   test --否--> output -> finish. The left return branch is the loop-back path;
%   there is no stepk -> output edge and no connection from the loop to output.
% solid segments: all box outlines and connector lines; dashed segments: none.
% labels: decision condition centered; “是” below the diamond and “否” above/right;
%   math labels centered in their boxes, with branch paths kept clear of text.
% directed edges (source): start->input, input->test; test->stepa (是),
% stepa->stepb->stepk->test (loop); test->output (否), output->finish.
\begin{tikzpicture}[
  >=Stealth,
  line width=0.45pt,
  line join=round,
  line cap=round,
  font=\normalsize,
  flowtext/.style={anchor=center, align=center,
                   execute at begin node={\setlength{\parindent}{0pt}}},
  flowbox/.style={flowtext, draw,
                  minimum height=0.60cm, inner xsep=4pt, inner ysep=2pt},
  process/.style={flowbox},
  startstop/.style={flowbox, rounded corners=0.18cm},
  io/.style={flowbox, trapezium, trapezium left angle=75, trapezium right angle=105,
             minimum height=0.70cm},
  decision/.style={flowbox, diamond, aspect=2.8, minimum width=2.75cm, minimum height=0.90cm},
]
  % Coordinates are a clean schematic layout; topology and labels follow the source figure.
  \node[startstop] (start) at (0,10.25) {开始};
  \node[io] (input) at (0,8.83) {输入 \(n=3,A=1,B=2,k=1\)};
  \node[decision]  (test)  at (0,7.40) {\(k\leq n\)};
  \node[process]   (stepa) at (0,5.97) {\(A=A+B\)};
  \node[process]   (stepb) at (0,4.54) {\(B=A+B\)};
  \node[process]   (stepk) at (0,3.11) {\(k=k+1\)};
  \node[io]        (output) at (0,1.68) {输出 \(B\)};
  \node[startstop] (finish) at (0,0.25) {结束};

  \draw[->] (start.south) -- (input.north);
  % The forward stem and loop return share a merge above the decision;
  % the short downward stem carries the sole arrow entering the diamond.
  \draw (input.south) -- (0,8.10);
  \draw[->] (0,8.10) -- (test.north);
  \draw[->] (test.south) -- node[flowtext, right] {是} (stepa.north);
  \draw[->] (stepa.south) -- (stepb.north);
  \draw[->] (stepb.south) -- (stepk.north);
  \draw[->] (output.south) -- (finish.north);

  % Loop-back branch from the increment box returns along the left side to the
  % upper test entry, matching the source PDF's horizontal arrow and short entry.
  \draw (stepk.west) -- (-1.80,3.11) -- (-1.80,8.10);
  \draw[->] (-1.80,8.10) -- (0,8.10);
  % False branch bypasses the loop and follows the right-side route into the
  % top of the output box; it then continues output -> finish.  There is no
  % k -> output connector.
  \draw[-] (test.east) -- node[flowtext, above] {否} (2.00,7.40)
    -- (2.00,2.38) -- (0,2.38);
  \draw[->] (0,2.38) -- (output.north);
\end{tikzpicture}
